\documentclass[addpoints]{exam}
% \documentclass[addpoints,12pt]{exam}

\usepackage[slovene]{babel}
\usepackage{amsfonts}
\usepackage{amssymb}
\usepackage{amsmath}
\usepackage{float}
\usepackage{graphicx}
\usepackage{enumitem}
\usepackage{color}
\usepackage{eurosym}


%Komentar naslednje vrstice glede na verzijo testa -- rešitve ali prostor za odgovore:
% \printanswers

% \linespread{2}


\pointsinrightmargin
\marginbonuspointname{}


\renewcommand{\solutiontitle}{\noindent\textbf{Rešitve in točkovanje:}\par\noindent}

\colorgrids
\definecolor{GridColor}{gray}{0.7}
\setlength{\gridsize}{7mm}

\extrawidth{5mm}
\setlength{\rightpointsmargin}{1.5cm}

\begin{document}

\runningfooter{}{}{}

\begin{center}
    \fbox{\fbox{\parbox{15cm}{\centering
    Popravljanje in pridobivanje ocen -- drugo pisno ocenjevanje znanja \\ 1. letnik -- test H \\ 17. junij 2026 \\ (Potence in izrazi; Deljivost)}}}
\end{center}

\vspace{0.1in}
\makebox[\textwidth]{Ime in priimek, oddelek:\enspace\hrulefill}


\begin{center}
    \hqword{Naloga}
    \hpword{Možne točke}
    % \chbpword{Dodatne točke}
    \hsword{Dosežene točke}
    \htword{Skupaj}  
    % \combinedpointtable[h][questions]
    \gradetable[h][questions]

\end{center}

\vspace{0.1in}


\begin{questions}

    % 1. vprašanje
    
    \question
    Izračunajte.
\begin{parts}
    \part[4]
    $-1^{2026}+\left((-2)^5+5^2-\left(7-3^2\right)^3\right)^2$

    \begin{solution}[5cm]
    $$\begin{aligned}
    \left(\left(-2\right)^3+\left(-5\right)^2\right)\cdot\left(-1\right)^{2025}-\left(\left(-3\right)^2\cdot\left(-2\right)-\left(-4\right)^2\right)-1^{2026}&=(-8+25)(-1)-(9(-2)-16)-1= \\ &=-17+34-1= \\ &=-18
    \end{aligned}$$
        Za pravilno uporabo pravila potenciranja z negativnim predznakom $1$ točka, pravilno potenciranje $1$ točka, pravilno množenje $1$ točka, končen rezultat $1$ točka.
    \end{solution}



    \part[3]
    $a^3(a^3-b^2)-a^6+(-a)^3b^2$

    \begin{solution}[4cm]
    $$5\cdot 2^{2n+1}-2\cdot 4^{n+1}-2\cdot 2^{2n+2}=5\cdot 2^{2n+1}-2^{2n+3}-2^{2n+3}=2^{2n+1}(5-2^2-2^2)=-3\cdot 2^{2n+1}$$
        Za pravilen zapis potenc števila $2$ $1$ točka, pravilno izpostavljanje skupnega faktorja $1$ točka, končen rezultat $1$ točka.
    \end{solution}

        
    \end{parts}


    % 2. vprašanje
    \question[5]
    Skrčite izraz in rezultat razstavite: $(3x-y)(3x+y)+2(2x-y)^2-5x(3x-2y)-y(4x+y)$.

    \begin{solution}[5cm]
        \begin{align*}
            (7a-2b)^2-(5a+2b)(5a-2b)-6(4a^2-b^2)&=49a^2-28ab+4b^2-25a^2+4b^2-24a^2+6b^2=\\
                &=14b^2-28ab= \\
                &=14b(b-2a)
        \end{align*}
        Za pravilno kvadriranje binoma $1$ točka, za pravilno uporabo pravila razlike kvadratov $1$ točka, pravilno množenje $1$ točka, pravilno izpostavljanje $1$ točka, pravilno uporabljena Vietova formula $1$ točka.
    \end{solution}
    

    % \newpage
    % 3. vprašanje
    \question
    Faktorizirajte izraze, kolikor je mogoče.
    \begin{parts}
        \part[2]
        $x^2-19x+90$

        \begin{solution}[4cm]
            $$2x^2-30x+108=2(x-6)(x-9)$$
            Za uporabo Vietovega pravila $1$ točka, pravilen rezultat $1$ točka.
        \end{solution}


        \part[2]
        $lk^2-64l$

        \begin{solution}[4cm]
            $$f^4-9g^2=(f^2-3g)(f^2+3g)$$
            Za pravilno izpostavljanje $1$ točka, pravilno uporabljena formula razlike kvadratov $1$ točka.
        \end{solution}


        \part[3]
        $j^3-3j^2+25j-75$

        \begin{solution}[4cm]
            $$c^3-4c^2-c+4=(c-1)(c+1)(c-4)$$
            Za uporabo pravila faktorizacije štiričlenika 2+2 $1$ točka, pravilna faktorizacija štiričlenika $1$ točka, pravilno razstavljena razlika kvadratov $1$ točka.
        \end{solution}


        \part[3]
        $64m-n^3m^4$

        \begin{solution}[4cm]
            $$m^5+27g^6m^2=m^2(m+3g^2)(m^2-6g^2+9g^4)$$
            Za pravilno izpostavljanje skupnega faktorja $1$ točka, uporaba formule razlike kubov $1$ točka, pravilen končen rezultat $1$ točka.
        \end{solution}


        \part[3]
        $x^2-6xy+9y^2-16$

        \begin{solution}[3cm]
            $$81-9b^2+12bc-4c^2=(9-3b+2c)(9+3b-2c)$$
            Za uporabo pravila faktorizacije štiričlenika 3+1 $1$ točka, pravilna uporaba formule $1$ točka, pravilno razstavljanje razlike kvadratov $1$ točka.
        \end{solution}


    \end{parts}


    \newpage

    % 4. vprašanje
    \question
    Določite neznano števko, da bo število $\overline{676s67s76s}$ deljivo
    \begin{parts}
        \part[2]
        s številom $3$.

        \begin{solution}[3.5cm]
            Za navedbo kriterija za deljivost z $2$ $1$ točka, uporaba kriterija in izračun $a\in\{0,2,4,6,8\}$ $1$ točka.
        \end{solution}


        \part[3]
        s številom $12$.

        \begin{solution}[5cm]
            Pravilno uporabljen kriterij za deljivost s $3$ $1$ točka, pravilno uporabljen kriterij za deljivost s $4$ $1$ točka, 
            pravilen zaključek in zapis kombinacij $(a,b)\in\left\{(2,0),(2,3),(2,6),(2,9),(6,2),(6,5),(6,8)\right\}$ $1$ točka.
        \end{solution}


    \end{parts}



% 5. vprašanje

    \question[4]
    Določite največji skupni delitelj $D$ in najmanjši skupni večkratnik $v$ števil $a=666$ in $b=1200$.
    Vrednost slednjega pustite zapisano v obliki produkta potenc prafaktorjev.

    \begin{solution}[6cm]
        Za zapisa $b=2\cdot 3^3\cdot 5^2\cdot 7$ in $c=2^2\cdot 3^2\cdot 5^4$ po $1$ točka. 
        Izračun $D(a,b,c)=3^2\cdot 5^2=225$ $1$ točka, in $v(a,b,c)=2^3\cdot 3^4\cdot 5^4\cdot 7\cdot 13$ $1$ točka.
    \end{solution}
    
    
% \newpage

    % 6. vprašanje
    \question[3]
    Število $x$ da ostnek $7$ pri deljenju z $8$, število $y$ pa ima pri deljenju s $4$ ostanek $1$. 
    Pokažite, da je število $x+y$ deljivo s številom $4$.

    
    \begin{solution}[4cm]

    \end{solution}

    % % 7. vprašanje
    % \question[4]
    %     Dokažite: Če kub poljubnega lihega naravnega števila zmanjšamo za $67$, bo dobljeno število deljivo z $2$.
        

    % \begin{solution}[3cm]

    % \end{solution}





\end{questions}



\end{document}